\chapter{Αξιολόγηση}
\label{ch:axiologisi}

Οι δύο εργασίες αναφοράς \cite{hadjichristofi2022, papias2023} αξιολογούν το σύστημά τους
κυρίως μέσω επίδειξης. Το παρόν κεφάλαιο προσθέτει τέσσερις ακόμη οπτικές: το κόστος των
λειτουργιών σε \en{gas}, δοκιμές αντοχής και κλιμάκωσης, μια προβολή του κόστους στο
\en{mainnet}, και την αντιστοίχιση των απαιτήσεων σε ελέγχους. Κλείνει με τους περιορισμούς της υλοποίησης.
Όλες οι μετρήσεις του κεφαλαίου έγιναν σε τοπικό δίκτυο, για τους λόγους που εξηγεί η
Ενότητα~\ref{sec:stress-method}· τα ποσά \en{gas} μεταφέρονται σε κάθε δίκτυο με απλή
αλλαγή τιμών, οι χρόνοι όχι.

\section{Κόστος συναλλαγών}
\label{sec:gas-eval}

\subsection{Μεθοδολογία}

Το κόστος μετρήθηκε σε τοπικό δίκτυο με ένα σενάριο που εκτελεί κάθε λειτουργία ως χωριστή
συναλλαγή και καταγράφει τη σύνοψη \en{gas} από τα αποτελέσματά της. Η τιμή αναφοράς του
\en{gas} ήταν 1.000 \en{MIST} ανά μονάδα. Όλα τα ποσά του κεφαλαίου δίνονται σε \en{SUI}
με έξι δεκαδικά ψηφία· 1 \en{SUI} $= 10^{9}$ \en{MIST}, οπότε το έκτο δεκαδικό ψηφίο
αντιστοιχεί σε 1.000 \en{MIST}. Για να είναι οι μετρήσεις καθαρές ελήφθησαν δύο
μέτρα. Το \en{faucet} παραδίδει το ποσό σε περισσότερα του ενός νομίσματα, και η πρώτη
συναλλαγή ενός λογαριασμού, αν αφεθεί στην προεπιλογή του \en{SDK}, τα συγχωνεύει όλα στο
νόμισμα που πληρώνει το \en{gas}· η διαγραφή των υπόλοιπων νομισμάτων παράγει επιστροφή
αποθήκευσης που δεν έχει σχέση με τη λειτουργία που μετράται. Γι' αυτό κάθε λογαριασμός
εκτέλεσε πρώτα μια ασήμαντη συναλλαγή εκτός μέτρησης, ώστε η συγχώνευση να γίνει εκεί, και
σε κάθε μετρούμενη συναλλαγή το \en{gas} πληρώθηκε ρητά από ένα μόνο, προκαθορισμένο
νόμισμα. Τα ποσά είναι ντετερμινιστικά: επανάληψη της
μέτρησης έδωσε πανομοιότυπες τιμές.

\subsection{Αποτελέσματα}

Ο Πίνακας~\ref{tab:gas} παρουσιάζει τις μετρήσεις. Οι τέσσερις στήλες του επεξηγούνται
αμέσως μετά.

\begin{table}[H]
\centering
\caption{Κόστος \en{gas} ανά λειτουργία, σε \en{SUI}}
\label{tab:gas}
\small
\begin{tabular}{lrrrr}
\toprule
Λειτουργία & Υπολογισμός & Αποθήκευση & Επιστροφή & Καθαρό \\
\midrule
\code{mint} $\times 1$ & 0,001000 & 0,003321 & 0,000978 & 0,003343 \\
\code{mint} $\times 3$ & 0,001000 & 0,007988 & 0,000978 & 0,008009 \\
\code{mint} $\times 10$ & 0,001000 & 0,024320 & 0,000978 & 0,024342 \\
\midrule
\code{place\_bid}: μόνο \en{SUI} & 0,001000 & 0,003914 & 0,000978 & 0,003936 \\
\code{place\_bid}: 1 \en{NFT} & 0,001000 & 0,009082 & 0,003288 & 0,006794 \\
\code{place\_bid}: 3 \en{NFTs} & 0,001000 & 0,019418 & 0,007908 & 0,012510 \\
\code{place\_bid}: 3 \en{NFTs} + \en{SUI} & 0,001000 & 0,019418 & 0,007908 & 0,012510 \\
\midrule
\code{accept\_bid}: μόνο \en{SUI} & 0,001000 & 0,004309 & 0,006185 & $-$0,000876 \\
\code{accept\_bid}: 1 \en{NFT} & 0,001000 & 0,005654 & 0,011301 & $-$0,004647 \\
\code{accept\_bid}: 3 \en{NFTs} & 0,001000 & 0,010321 & 0,021534 & $-$0,010213 \\
\code{accept\_bid}: 3 \en{NFTs} + \en{SUI} & 0,001000 & 0,011309 & 0,021534 & $-$0,009225 \\
\midrule
\code{reject\_bid}: 1 \en{NFT} & 0,001000 & 0,005654 & 0,011301 & $-$0,004647 \\
\code{cancel\_bid}: 1 \en{NFT} & 0,001000 & 0,003321 & 0,008991 & $-$0,004670 \\
\bottomrule
\end{tabular}
\end{table}

Οι στήλες του πίνακα αντιστοιχούν στα σκέλη του μοντέλου \en{gas} της
Ενότητας~\ref{sec:gas}:

\begin{description}[leftmargin=1.2em,itemsep=2pt]
  \item[Υπολογισμός.] Το κόστος εκτέλεσης της συναλλαγής: οι μονάδες υπολογισμού που
    κατανάλωσαν οι εντολές της, επί την τιμή αναφοράς του \en{gas} της εποχής. Στο τοπικό
    δίκτυο η τιμή αυτή είναι 1.000 \en{MIST} ανά μονάδα, και όλες οι λειτουργίες του
    συμβολαίου πέφτουν στο κατώτατο κλιμάκιο των 1.000 μονάδων: $1.000 \times 1.000$
    \en{MIST} $= 0{,}001$ \en{SUI}.
  \item[Αποθήκευση.] Το ποσό που δεσμεύεται ως εγγύηση για τα \en{bytes} που γράφει η
    συναλλαγή στην αλυσίδα: κάθε αντικείμενο που δημιουργείται, καθώς και κάθε υπάρχον
    αντικείμενο που τροποποιείται, συμπεριλαμβανομένου του νομίσματος που πληρώνει το
    \en{gas}. Χρεώνεται με σταθερή τιμή ανά \en{byte}.
  \item[Επιστροφή.] Το μέρος της εγγύησης αποθήκευσης που επιστρέφεται στον αποστολέα
    επειδή η συναλλαγή διαγράφει ή ξαναγράφει αντικείμενα για τα οποία είχε καταβληθεί
    εγγύηση στο παρελθόν, από οποιονδήποτε. Το 99\% της εγγύησης επιστρέφεται· το 1\%
    παρακρατείται. Η επιστροφή πιστώνεται σε όποιον υποβάλλει τη συναλλαγή της
    διαγραφής, όχι κατ' ανάγκην σε όποιον είχε πληρώσει την εγγύηση.
  \item[Καθαρό.] Το ποσό που αφαιρέθηκε τελικά από το νόμισμα \en{gas} του αποστολέα:
    υπολογισμός συν αποθήκευση μείον επιστροφή. Αρνητική τιμή σημαίνει ότι ο αποστολέας
    εισέπραξε περισσότερα από όσα πλήρωσε, επειδή απελευθέρωσε εγγύηση που είχε καταβάλει
    κάποιος άλλος· αυτό συμβαίνει σε όλες τις λειτουργίες που κλείνουν μια προσφορά.
\end{description}

\subsection{Ανάλυση}

\paragraph{Ο υπολογισμός είναι σταθερός και αμελητέος.} Όλες οι λειτουργίες χρεώνονται το
ελάχιστο κόστος υπολογισμού, 0{,}001 \en{SUI}, ανεξάρτητα από το πλήθος των \en{NFTs}. Η
λογική του συμβολαίου είναι απλή, και ακόμη και οι βρόχοι πάνω από τρία ή δέκα αντικείμενα
δεν ξεπερνούν το κατώτατο κλιμάκιο χρέωσης.

\paragraph{Το κόστος καθορίζεται από την αποθήκευση.} Από τις διαφορές μεταξύ των γραμμών
του πίνακα προκύπτει το κόστος αποθήκευσης κάθε είδους αντικειμένου
(Πίνακας~\ref{tab:storage-units}). Κάθε γραμμή του Πίνακα~\ref{tab:gas} αναλύεται ακριβώς
σε άθροισμα τέτοιων όρων. Για παράδειγμα, η υποβολή προσφοράς με ένα \en{NFT} αποθηκεύει
το \code{Bid}, ένα δυναμικό πεδίο, το \en{NFT} με νέο ιδιοκτήτη και το νόμισμα του \en{gas}:
$0{,}002926 + 0{,}002835 + 0{,}002333 + 0{,}000988 = 0{,}009082$ \en{SUI}.

\begin{table}[H]
\centering
\caption{Κόστος αποθήκευσης ανά αντικείμενο, όπως προκύπτει από τις μετρήσεις}
\label{tab:storage-units}
\small
\begin{tabular}{lr}
\toprule
Αντικείμενο & Αποθήκευση (\en{SUI}) \\
\midrule
\code{MarketplaceNFT} (με διεύθυνση εικόνας 80 χαρακτήρων) & 0,002333 \\
\code{Bid} & 0,002926 \\
Δυναμικό πεδίο αντικειμένου (ανά δεσμευμένο \en{NFT}) & 0,002835 \\
\code{Coin<SUI>} & 0,000988 \\
\bottomrule
\end{tabular}
\end{table}

Το μέγεθος του \code{Bid} αξίζει μια σύγκριση. Σε προγενέστερη έκδοση του συμβολαίου η
δομή περιείχε επιπλέον τα ονόματα των δύο τύπων ως συμβολοσειρές και είχε την ικανότητα
\code{store}, και η αποθήκευσή της κόστιζε 0{,}004294 \en{SUI}. Η αφαίρεση των δύο
πεδίων, που επαναλάμβαναν πληροφορία ήδη διαθέσιμη στον τύπο του αντικειμένου
(Ενότητα~\ref{sec:data-model}), μείωσε το κόστος κατά 32\%. Η αποθήκευση στο \en{Sui}
χρεώνεται ανά \en{byte}, και κάθε πλεονάζον πεδίο έχει τίμημα.

\paragraph{Η δέσμευση είναι εγγύηση, όχι δαπάνη.} Το κόστος αποθήκευσης επιστρέφεται κατά
99\% όταν το αντικείμενο διαγραφεί (Ενότητα~\ref{sec:gas}). Επειδή κάθε προσφορά
καταλήγει σε διαγραφή του \code{Bid} και των πεδίων του, το μεγαλύτερο μέρος όσων
καταβάλλονται κατά την υποβολή ανακτάται κατά τον τερματισμό. Ο
Πίνακας~\ref{tab:lifecycle-cost} δείχνει το συνολικό κόστος ενός πλήρους κύκλου ζωής για
προσφορά με ένα \en{NFT}.

\begin{table}[H]
\centering
\caption{Καθαρό κόστος ανά κύκλο ζωής προσφοράς με ένα \en{NFT}, σε \en{SUI}}
\label{tab:lifecycle-cost}
\small
\begin{tabular}{lrrr}
\toprule
Κύκλος ζωής & Προσφέρων & Κάτοχος & Σύνολο \\
\midrule
Υποβολή και αποδοχή & 0,006794 & $-$0,004647 & 0,002147 \\
Υποβολή και απόρριψη & 0,006794 & $-$0,004647 & 0,002147 \\
Υποβολή και ακύρωση & 0,002124 & 0 & 0,002124 \\
\bottomrule
\end{tabular}
\end{table}

Μια πλήρης ανταλλαγή κοστίζει συνολικά περίπου 0{,}002 \en{SUI}: δύο φορές το ελάχιστο
κόστος υπολογισμού και το μη επιστρεφόμενο 1\% της αποθήκευσης. Το κόστος είναι
ανεξάρτητο από την αξία των στοιχείων που ανταλλάσσονται, και το πρωτόκολλο δεν εισπράττει
τίποτε (ΜΛΑ4).

\paragraph{Η εγγύηση επιστρέφεται σε όποιον κλείνει την προσφορά.} Ο
Πίνακας~\ref{tab:lifecycle-cost} αναδεικνύει μια ασυμμετρία που δεν είχε προβλεφθεί στον
σχεδιασμό. Την εγγύηση αποθήκευσης την καταβάλλει ο προσφέρων, αλλά η επιστροφή πηγαίνει
στον αποστολέα της συναλλαγής που διαγράφει τα αντικείμενα. Στην ακύρωση αυτός είναι ο
ίδιος ο προσφέρων. Στην αποδοχή και στην απόρριψη είναι ο κάτοχος, ο οποίος εισπράττει
περίπου 0{,}005 \en{SUI} που είχε καταβάλει ο προσφέρων. Το ποσό είναι μικρό και η
συνέπεια έχει και θετική όψη: ο κάτοχος έχει ελαφρύ κίνητρο να απορρίπτει ρητά τις
προσφορές που δεν τον ενδιαφέρουν, αποδεσμεύοντας έτσι τα στοιχεία του προσφέροντος.
Παραμένει ωστόσο μια μεταφορά αξίας που οι χρήστες δεν βλέπουν, και θα έπρεπε είτε να
εμφανίζεται στη διεπαφή είτε να αντισταθμίζεται από το συμβόλαιο.

\paragraph{Κλιμάκωση με το πλήθος των \en{NFTs}.} Κάθε επιπλέον δεσμευμένο \en{NFT}
προσθέτει στην υποβολή περίπου 0{,}0029 \en{SUI} καθαρά, σχεδόν εξ ολοκλήρου ως
εγγύηση για το δυναμικό πεδίο. Η σχέση είναι γραμμική, και δεν παρατηρήθηκε αύξηση του
κόστους υπολογισμού έως τα δέκα αντικείμενα ανά συναλλαγή.

\subsection{Χρόνος απόκρισης}

Στο τοπικό δίκτυο, ο χρόνος από την υποβολή μιας συναλλαγής έως την επιστροφή του
αποτελέσματος ήταν περίπου 0{,}20 δευτερόλεπτα, και έως ότου η νέα κατάσταση γίνει ορατή
στα ερωτήματα περίπου 0{,}50 δευτερόλεπτα, με μικρή διακύμανση μεταξύ λειτουργιών. Οι
τιμές αυτές αφορούν δίκτυο που εκτελείται σε έναν υπολογιστή, με έναν επικυρωτή, και δεν
αντιπροσωπεύουν δημόσιο δίκτυο (Ενότητα~\ref{sec:sui-networks}). Δείχνουν πάντως ότι η εφαρμογή δεν προσθέτει δική της καθυστέρηση: ο
χρόνος που βιώνει ο χρήστης καθορίζεται από το πορτοφόλι και το δίκτυο.


\section{Δοκιμές αντοχής και κλιμάκωσης}
\label{sec:stress}

Οι μετρήσεις της προηγούμενης ενότητας αφορούν τη συνήθη χρήση: μία προσφορά, λίγα
\en{NFTs}, ένας χρήστης τη φορά. Η παρούσα ενότητα πιέζει το σύστημα σε τέσσερις
διαστάσεις, το μέγεθος μιας προσφοράς, το μέγεθος μιας παρτίδας δημιουργίας, το πλήθος
ταυτόχρονων χρηστών και το μέγεθος του ιστορικού, και καταγράφει πού αλλάζει η
συμπεριφορά του και πού βρίσκονται τα όριά του. Όλες οι μετρήσεις έγιναν σε τοπικό
δίκτυο· η Ενότητα~\ref{sec:mainnet-projection} τις προβάλλει στο \en{mainnet}.

\subsection{Μεθοδολογία και επιλογή δικτύου}
\label{sec:stress-method}

Οι δοκιμές εκτελέσθηκαν από ένα σενάριο \en{TypeScript} πάνω στο \en{SDK}, στο ίδιο
τοπικό δίκτυο ενός επικυρωτή της Ενότητας~\ref{sec:gas-eval} (\en{sui} 1.67.1,
πρωτόκολλο 114), με ένα μόνο πακέτο δημοσιευμένο και χωρίς άλλο φορτίο. Ένας πλούσιος
λογαριασμός, χρηματοδοτημένος από το τοπικό \en{faucet}, μοίρασε νομίσματα σε όσους
λογαριασμούς χρειάστηκε κάθε δοκιμή, έως 250 ταυτόχρονα, με ένα \en{PTB} ανά 100
παραλήπτες. Κάθε λογαριασμός υπογράφει τις δικές του συναλλαγές με δικό του νόμισμα
\en{gas}, ώστε οι ταυτόχρονες συναλλαγές να μην συγκρούονται μεταξύ τους παρά μόνο εκεί
που η δοκιμή το επιδιώκει. Οι χρόνοι μετρούνται στον πελάτη, από την αποστολή έως την
επιστροφή των αποτελεσμάτων.

Η επιλογή του τοπικού δικτύου αντί του \en{testnet} ήταν συνειδητή, για τέσσερις λόγους.
Πρώτον, η χρηματοδότηση: οι δοκιμές χρειάσθηκαν εκατοντάδες λογαριασμούς και περισσότερες
από δύο χιλιάδες συναλλαγές, ενώ το \en{faucet} του \en{testnet} διαθέτει μικρό ποσό ανά
αίτημα με όριο ρυθμού ανά διεύθυνση δικτύου (Ενότητα~\ref{sec:sui-networks}). Δεύτερον,
οι δημόσιοι πλήρεις κόμβοι περιορίζουν τον ρυθμό αιτημάτων, οπότε μια ριπή εκατό
ταυτόχρονων συναλλαγών θα μετρούσε τον περιορισμό του κόμβου και όχι το σύστημα. Τρίτον,
η επαναληψιμότητα: το τοπικό δίκτυο ξεκινά από καθαρή κατάσταση και δεν μοιράζεται το
φορτίο του με κανέναν, ώστε δύο εκτελέσεις να δίνουν συγκρίσιμους χρόνους. Τέταρτον, το
κόστος σε \en{gas} δεν εξαρτάται από το δίκτυο: οι μονάδες υπολογισμού και τα \en{bytes}
αποθήκευσης είναι ίδια παντού για την ίδια συναλλαγή, και μόνο οι τιμές ανά μονάδα
διαφέρουν. Έτσι οι μετρήσεις κόστους μεταφέρονται με ακρίβεια σε οποιοδήποτε δίκτυο,
όπως γίνεται στην Ενότητα~\ref{sec:mainnet-projection}, ενώ οι χρόνοι και ο ρυθμός
διεκπεραίωσης αφορούν μόνο το περιβάλλον που μετρήθηκε. Το τίμημα είναι ότι η εικόνα
για την καθυστέρηση ενός δημόσιου δικτύου μπορεί να είναι μόνο προσεγγιστική.

\subsection{Μέγεθος προσφοράς}
\label{sec:stress-bidsize}

Η πρώτη δοκιμή υποβάλλει και αποδέχεται προσφορές με 1 έως 500 δεσμευμένα \en{NFTs}
(Πίνακας~\ref{tab:stress-bidsize}, Σχήμα~\ref{fig:stress-bidsize}). Και οι 500
δεσμεύσεις πέρασαν· τα όρια του πρωτοκόλλου που θα σταματούσαν την κλιμάκωση είναι τα
512 ορίσματα ανά εντολή και τα 2.048 αντικείμενα εισόδου ανά συναλλαγή, και δεν
προσεγγίσθηκαν διαφορετικά.

\begin{table}[H]
\centering
\caption{Κόστος και χρόνος προσφοράς με $n$ δεσμευμένα \en{NFTs}, σε \en{SUI} (τοπικό δίκτυο)}
\label{tab:stress-bidsize}
\small
\begin{tabular}{rrrrrrr}
\toprule
& \multicolumn{3}{c}{\code{place\_bid}} & \multicolumn{3}{c}{\code{accept\_bid}} \\
\cmidrule(lr){2-4}\cmidrule(lr){5-7}
$n$ & Υπολογισμός & Καθαρό & ms & Υπολογισμός & Καθαρό & ms \\
\midrule
1 & 0,001000 & 0,006794 & 156 & 0,001000 & $-$0,004647 & 161 \\
5 & 0,001000 & 0,018226 & 210 & 0,001000 & $-$0,015780 & 213 \\
10 & 0,001000 & 0,032516 & 217 & 0,001000 & $-$0,029696 & 209 \\
25 & 0,001000 & 0,075386 & 210 & 0,001000 & $-$0,071444 & 209 \\
50 & 0,001000 & 0,146838 & 213 & 0,001000 & $-$0,141024 & 212 \\
100 & 0,001410 & 0,290151 & 208 & 0,001000 & $-$0,280184 & 214 \\
250 & 0,065000 & 0,782460 & 204 & 0,012900 & $-$0,685752 & 202 \\
500 & 0,182400 & 1,614393 & 293 & 0,055700 & $-$1,338732 & 276 \\
\bottomrule
\end{tabular}
\end{table}

\begin{figure}[H]
\centering
\begin{tikzpicture}
\begin{loglogaxis}[width=0.8\textwidth,height=6.2cm,xlabel={δεσμευμένα \en{NFTs} ανά προσφορά},ylabel={\en{SUI}},legend pos=north west,legend style={font=\footnotesize},grid=major,xtick={1,5,10,25,50,100,250,500},xticklabels={1,5,10,25,50,100,250,500},tick label style={font=\footnotesize},label style={font=\footnotesize}]
\addplot+[mark=*] coordinates {(1,0.00679386) (5,0.01822578) (10,0.03251568) (25,0.075386444) (50,0.146837844) (100,0.290150644) (250,0.782460292) (500,1.614393444)};
\addplot+[mark=square*] coordinates {(1,0.001) (5,0.001) (10,0.001) (25,0.001) (50,0.001) (100,0.00141) (250,0.065) (500,0.1824)};
\legend{καθαρό κόστος \code{place\_bid}, υπολογισμός \code{place\_bid}}
\end{loglogaxis}
\end{tikzpicture}
\caption{Κόστος υποβολής προσφοράς ως συνάρτηση του πλήθους δεσμευμένων \en{NFTs}
(λογαριθμικοί άξονες). Το καθαρό κόστος είναι γραμμικό· ο υπολογισμός μένει στο ελάχιστο
κλιμάκιο έως τα 50 και εκτοξεύεται πάνω από τα 100.}
\label{fig:stress-bidsize}
\end{figure}

Τρία πράγματα φαίνονται. Το \emph{καθαρό} κόστος είναι γραμμικό στο $n$, περίπου 0{,}0029
\en{SUI} ανά \en{NFT}, και σχεδόν εξ ολοκλήρου εγγύηση αποθήκευσης που
επιστρέφεται κατά την αποδοχή: η αποδοχή προσφοράς 500 \en{NFTs} επιστρέφει 1{,}34
\en{SUI} στον αποδέκτη. Ο \emph{υπολογισμός} όμως δεν είναι γραμμικός. Μένει στο
ελάχιστο κλιμάκιο των 0{,}001 \en{SUI} έως τα 50 \en{NFTs}, ανεβαίνει ελαφρά
στα 100, και στα 250 και 500 γίνεται 65 και 182 φορές το ελάχιστο: η χρέωση ανά
\en{byte} για την ανάγνωση και εγγραφή εκατοντάδων αντικειμένων και την προσάρτηση
ισάριθμων δυναμικών πεδίων κυριαρχεί. Η αποδοχή είναι τρεις φορές φθηνότερη σε
υπολογισμό από την υποβολή στο ίδιο $n$, επειδή αφαιρεί πεδία αντί να τα δημιουργεί.
Τέλος, ο χρόνος εκτέλεσης δεν αλλάζει ουσιαστικά, από 0{,}16 έως 0{,}29 δευτερόλεπτα· σε
τοπικό δίκτυο η καθυστέρηση καθορίζεται από τη συναίνεση, όχι από το μέγεθος της
συναλλαγής. Για τη χρήση που προβλέπει η εφαρμογή, λίγα \en{NFTs} ανά προσφορά, το
σύστημα βρίσκεται βαθιά μέσα στην επίπεδη περιοχή.

\subsection{Παρτίδες δημιουργίας}

Η δεύτερη δοκιμή δημιουργεί 1 έως 200 \en{NFTs} σε ένα \en{PTB}, με μία εντολή
\code{mint} ανά αντικείμενο και μία κοινή \code{TransferObjects}
(Πίνακας~\ref{tab:stress-mint}). Ο υπολογισμός μένει στο ελάχιστο κλιμάκιο ακόμη και για
200 κλήσεις, και το κόστος ανά \en{NFT} συγκλίνει στα 0{,}00233 \en{SUI}, την
αποθήκευση του ίδιου του αντικειμένου· το σταθερό κόστος της συναλλαγής μοιράζεται σε
όλο και περισσότερα αντικείμενα. Ο χρόνος παραμένει κάτω από 0{,}2 δευτερόλεπτα. Το
όριο εδώ είναι οι 1.024 εντολές ανά \en{PTB} και τα 512 ορίσματα της τελικής μεταβίβασης.

\begin{table}[H]
\centering
\caption{Δημιουργία $n$ \en{NFTs} σε μία συναλλαγή, σε \en{SUI} (τοπικό δίκτυο)}
\label{tab:stress-mint}
\small
\begin{tabular}{rrrrr}
\toprule
$n$ & Υπολογισμός & Καθαρό & Καθαρό ανά \en{NFT} & ms \\
\midrule
1 & 0,001000 & 0,003328 & 0,003328 & 128 \\
10 & 0,001000 & 0,024190 & 0,002419 & 162 \\
50 & 0,001000 & 0,117214 & 0,002344 & 172 \\
100 & 0,001000 & 0,233494 & 0,002335 & 151 \\
200 & 0,001000 & 0,466814 & 0,002334 & 174 \\
\bottomrule
\end{tabular}
\end{table}

\subsection{Ταυτόχρονοι χρήστες}
\label{sec:stress-throughput}

Η τρίτη δοκιμή προσομοιώνει αγορά με φορτίο: 100 προσφέροντες υποβάλλουν ταυτόχρονα από
μία προσφορά ο καθένας προς 100 διαφορετικούς κατόχους, και έπειτα οι 100 κάτοχοι
αποδέχονται ταυτόχρονα. Και οι 200 συναλλαγές ολοκληρώθηκαν επιτυχώς
(Πίνακας~\ref{tab:stress-throughput}). Κάθε ριπή των 100 χρειάσθηκε κάτω από ένα
δευτερόλεπτο συνολικά, δηλαδή περίπου 104 συναλλαγές ανά δευτερόλεπτο από έναν μόνο
πελάτη, και ο διάμεσος χρόνος ανά συναλλαγή μέσα στη ριπή ήταν 0{,95} δευτερόλεπτα, σχεδόν
όσο και η ριπή ολόκληρη: ο κόμβος παραλαμβάνει τις 100 συναλλαγές σχεδόν μαζί και τις
διεκπεραιώνει μαζί, όχι σειριακά. Το μέτρο αυτό αφορά τον ρυθμό που μπορεί να παραγάγει
ένας πελάτης σε έναν τοπικό κόμβο, και δεν είναι μέτρηση της χωρητικότητας του δικτύου.
Το ουσιώδες εύρημα είναι η απουσία αποτυχιών: επειδή κάθε προσφορά είναι χωριστό
κοινόχρηστο αντικείμενο και κάθε αποστολέας έχει δικό του νόμισμα, εκατό ανεξάρτητες
ανταλλαγές δεν διεκδικούν κανέναν κοινό πόρο.

\begin{table}[H]
\centering
\caption{Εκατό ταυτόχρονες συναλλαγές ανά ροή (τοπικό δίκτυο)}
\label{tab:stress-throughput}
\footnotesize
\begin{tabular}{lrrrrr}
\toprule
Ροή & Επιτυχείς & Αποτυχίες & Σύνολο & Συν./s & Διάμεσος \\
\midrule
100 $\times$ \code{place\_bid} & 100 & 0 & 0{,}96 s & 104 & 0{,}95 s \\
100 $\times$ \code{accept\_bid} & 100 & 0 & 0{,}96 s & 104 & 0{,}95 s \\
\bottomrule
\end{tabular}
\end{table}

\subsection{Ανταγωνισμός για την ίδια προσφορά}
\label{sec:stress-contention}

Η τέταρτη δοκιμή είναι η μόνη που επιδιώκει σύγκρουση: $k$ λογαριασμοί, από 2 έως 20,
προσπαθούν ταυτόχρονα να αποδεχθούν την ίδια ανοικτή προσφορά, σε τρεις γύρους για κάθε
$k$ (Πίνακας~\ref{tab:stress-contention}). Σε κάθε έναν από τους δώδεκα γύρους ακριβώς
ένας πέτυχε και όλοι οι υπόλοιποι απέτυχαν με τον ίδιο λόγο, \code{InputObjectDeleted}:
η συναίνεση διέταξε τις $k$ συναλλαγές, η πρώτη διέγραψε το \code{Bid}, και οι επόμενες
βρήκαν την είσοδό τους να μην υπάρχει. Καμία αποτυχημένη δεν άφησε ίχνος πέραν του
\en{gas} της. Αυτή είναι η πειραματική επιβεβαίωση του επιχειρήματος της
Ενότητας~\ref{sec:security} για την ταυτόχρονη πρόσβαση, και δείχνει ότι το πρωτόκολλο
δεν χρειάζεται κανένα δικό του κλείδωμα: το παρέχει το μοντέλο αντικειμένων.

\begin{table}[H]
\centering
\caption{$k$ ταυτόχρονες αποδοχές της ίδιας προσφοράς, τρεις γύροι ανά $k$}
\label{tab:stress-contention}
\small
\begin{tabular}{rrrrl}
\toprule
$k$ & Γύροι & Επιτυχείς ανά γύρο & Αποτυχίες συνολικά & Αιτία αποτυχίας \\
\midrule
2 & 3 & 1 & 3 & \code{InputObjectDeleted} \\
5 & 3 & 1 & 12 & \code{InputObjectDeleted} \\
10 & 3 & 1 & 27 & \code{InputObjectDeleted} \\
20 & 3 & 1 & 57 & \code{InputObjectDeleted} \\
\bottomrule
\end{tabular}
\end{table}

\subsection{Κόστος ανακάλυψης με το ιστορικό}
\label{sec:stress-discovery}

Η τελευταία δοκιμή δεν πιέζει το συμβόλαιο αλλά την εφαρμογή, και ποσοτικοποιεί τον
περιορισμό της ανακάλυψης μέσω συμβάντων (Ενότητα~\ref{sec:reading}). Μετρήθηκε ο χρόνος
του ερωτήματος εισερχόμενων προσφορών ενός κατόχου, όπως το εκτελεί η εφαρμογή
(ανάγνωση όλων των \code{BidPlaced}, ανάκτηση των αντικειμένων, ανάκτηση των στόχων),
ενώ πενήντα λογαριασμοί γέμιζαν την αλυσίδα με ενεργές προσφορές
(Πίνακας~\ref{tab:stress-discovery}, Σχήμα~\ref{fig:stress-discovery}).

\begin{table}[H]
\centering
\caption{Χρόνος του ερωτήματος εισερχόμενων προσφορών ως συνάρτηση του ιστορικού}
\label{tab:stress-discovery}
\small
\begin{tabular}{rrr}
\toprule
Συμβάντα \code{BidPlaced} & Ενεργές προσφορές & Διάμεσος χρόνος (3 επαναλήψεις) \\
\midrule
237 & 4 & 14 ms \\
487 & 254 & 32 ms \\
787 & 554 & 47 ms \\
1.487 & 1.254 & 87 ms \\
\bottomrule
\end{tabular}
\end{table}

\begin{figure}[H]
\centering
\begin{tikzpicture}
\begin{axis}[width=0.72\textwidth,height=5.4cm,xlabel={ενεργές προσφορές στην αλυσίδα},ylabel={ms},grid=major,ymin=0,tick label style={font=\footnotesize},label style={font=\footnotesize}]
\addplot+[mark=*] coordinates {(4,14) (254,32) (554,47) (1254,87)};
\end{axis}
\end{tikzpicture}
\caption{Ο χρόνος του ερωτήματος εισερχόμενων προσφορών αυξάνει γραμμικά με το πλήθος
των προσφορών στην αλυσίδα, περίπου 58 ms ανά χίλιες προσφορές σε τοπικό κόμβο.}
\label{fig:stress-discovery}
\end{figure}

Η σχέση είναι γραμμική, περίπου 58 χιλιοστά του δευτερολέπτου ανά χίλιες ενεργές
προσφορές, και ο χρόνος παραμένει κάτω από το ένα δέκατο του δευτερολέπτου στις 1.250. Ο
κόμβος όμως είναι τοπικός· σε δημόσιο κόμβο κάθε σελίδα συμβάντων και κάθε παρτίδα
αντικειμένων είναι ένα αίτημα δικτύου, και το ίδιο ερώτημα θα κόστιζε δεκάδες τέτοια
αιτήματα. Η καμπύλη επιβεβαιώνει αυτό που η Ενότητα~\ref{sec:limitations} διατυπώνει
ποιοτικά: η προσέγγιση αρκεί για επίδειξη και μικρή κοινότητα, και χρειάζεται
\en{indexer} πέρα από μερικές χιλιάδες προσφορές.

\section{Προβολή στο mainnet}
\label{sec:mainnet-projection}

Το σύστημα δεν δοκιμάσθηκε στο \en{mainnet}, και δεν θα έπρεπε (Ενότητα~\ref{sec:dev-flow}).
Το κόστος του εκεί μπορεί ωστόσο να υπολογισθεί με ακρίβεια, επειδή το \en{gas} του
\en{Sui} είναι γινόμενο μετρημένων ποσοτήτων επί δημοσιευμένες τιμές
(Ενότητα~\ref{sec:gas}): οι μονάδες υπολογισμού και τα \en{bytes} αποθήκευσης κάθε
συναλλαγής προκύπτουν από τις μετρήσεις, και οι τιμές ανά μονάδα διαβάζονται από το ίδιο
το δίκτυο. Ο Πίνακας~\ref{tab:net-params} δείχνει τις παραμέτρους των τριών δικτύων, όπως
διαβάστηκαν από τους πλήρεις κόμβους τους στις 16 Οκτωβρίου 2026.

\begin{table}[H]
\centering
\caption{Παράμετροι \en{gas} ανά δίκτυο (16 Οκτωβρίου 2026)}
\label{tab:net-params}
\small
\begin{tabular}{lrrr}
\toprule
Παράμετρος & \en{Localnet} (μετρήσεις) & \en{Testnet} & \en{Mainnet} \\
\midrule
Έκδοση πρωτοκόλλου & 114 & 139 & 138 \\
Τιμή αναφοράς \en{gas} (\en{MIST}/μονάδα) & 1.000 & 1.000 & 100 \\
Τιμή αποθήκευσης (\en{MIST}/μονάδα) & 76 & 76 & 76 \\
Ποσοστό επιστροφής αποθήκευσης & 99\% & 99{,}99\% & 99{,}99\% \\
\bottomrule
\end{tabular}
\end{table}

Δύο διαφορές έχουν σημασία. Η τιμή αναφοράς του \en{mainnet} είναι δέκα φορές
χαμηλότερη από του τοπικού δικτύου, οπότε ο υπολογισμός κοστίζει το ένα δέκατο. Και το
ποσοστό επιστροφής αποθήκευσης έχει αυξηθεί στις νεότερες εκδόσεις του πρωτοκόλλου από
99\% σε 99{,}99\%, οπότε το μη επιστρεφόμενο μέρος της εγγύησης, που στο τοπικό δίκτυο
ήταν το 1\%, γίνεται το ένα εκατοστό του. Η τιμή αποθήκευσης είναι ίδια. Η προβολή
υπολογίζει λοιπόν, για κάθε μετρημένη συναλλαγή, υπολογισμό ίσο με τις μετρημένες μονάδες
επί 100, αποθήκευση αμετάβλητη, και επιστροφή ίση με το 99{,}99\% της αποθήκευσης που
είχαν καταβάλει τα αντικείμενα που διαγράφονται ή τροποποιούνται. Ο
Πίνακας~\ref{tab:mainnet-cost} δίνει το αποτέλεσμα, μαζί με το ισοδύναμο σε δολάρια με
την ισοτιμία της ίδιας ημέρας~\cite{coingecko-sui}.

\begin{table}[H]
\centering
\caption{Προβολή του καθαρού κόστους ανά λειτουργία στο \en{mainnet}}
\label{tab:mainnet-cost}
\footnotesize
\begin{tabular}{lrrr}
\toprule
Λειτουργία & Τοπικό (\en{SUI}) & \en{Mainnet} (\en{SUI}) & \en{Mainnet} (USD) \\
\midrule
\code{mint} $\times 1$ & 0,003343 & 0,002433 & 0{,}0025 \\
\code{place\_bid}: μόνο \en{SUI} & 0,003936 & 0,003026 & 0{,}0031 \\
\code{place\_bid}: 1 \en{NFT} & 0,006794 & 0,005861 & 0{,}0061 \\
\code{place\_bid}: 3 \en{NFTs} & 0,012510 & 0,011531 & 0{,}0120 \\
\code{accept\_bid}: 1 \en{NFT} & $-$0,004647 & $-$0,005660 & $-$0{,}0059 \\
\code{accept\_bid}: 3 \en{NFTs} & $-$0,010213 & $-$0,011328 & $-$0{,}0118 \\
\code{reject\_bid}: 1 \en{NFT} & $-$0,004647 & $-$0,005660 & $-$0{,}0059 \\
\code{cancel\_bid}: 1 \en{NFT} & $-$0,004670 & $-$0,005660 & $-$0{,}0059 \\
\midrule
Πλήρης κύκλος: υποβολή και αποδοχή & 0,002147 & 0,000201 & 0{,}0002 \\
Πλήρης κύκλος: υποβολή και ακύρωση & 0,002124 & 0,000201 & 0{,}0002 \\
\bottomrule
\end{tabular}
\end{table}

Η εικόνα που προκύπτει είναι ακόμη ευνοϊκότερη από αυτή του τοπικού δικτύου. Μια πλήρης
ανταλλαγή με ένα \en{NFT} κοστίζει στο \en{mainnet} περίπου 0{,}0002 \en{SUI}, δηλαδή
δύο εκατοστά του σεντ: σχεδόν εξ ολοκλήρου οι δύο ελάχιστες χρεώσεις
υπολογισμού των 0{,}0001 \en{SUI}, αφού η εγγύηση αποθήκευσης επιστρέφεται πλέον στο
σύνολό της πλην αμελητέου υπολοίπου. Η ασυμμετρία της Ενότητας~\ref{sec:gas-eval}
παραμένει και μεγαλώνει ελαφρά: ο κάτοχος που αποδέχεται ή απορρίπτει εισπράττει περίπου
0{,}0057 \en{SUI} (0{,}006 δολάρια) που είχε καταβάλει ο προσφέρων. Για τις
μεγάλες προσφορές της Ενότητας~\ref{sec:stress-bidsize}, η υποβολή με 100 \en{NFTs}
δεσμεύει προσωρινά 0{,}29 \en{SUI} (0{,}30 δολάρια) και με 500 \en{NFTs} 1{,}44 \en{SUI}
(1{,}49 δολάρια), ποσά που επιστρέφονται σχεδόν ολόκληρα κατά τον τερματισμό· το μόνιμο
κόστος των 500 \en{NFTs} είναι ο υπολογισμός, 0{,}0182 \en{SUI} στην υποβολή
και 0{,}0056 στην αποδοχή.

Για τον ρυθμό διεκπεραίωσης και την καθυστέρηση δεν υπάρχει αντίστοιχος ακριβής
υπολογισμός, και η εκτίμηση πρέπει να είναι ποιοτική. Οι 100 ταυτόχρονες ανταλλαγές της
Ενότητας~\ref{sec:stress-throughput} είναι αμελητέο φορτίο για το δίκτυο, που σχεδιάσθηκε
για τάξεις μεγέθους περισσότερες συναλλαγές ανά δευτερόλεπτο~\cite{sui-whitepaper,
mysticeti2023}· ο περιορισμός σε πραγματική χρήση θα ήταν ο ρυθμός αιτημάτων που δέχεται
ο δημόσιος κόμβος από έναν πελάτη. Κάθε συναλλαγή του πρωτοκόλλου αγγίζει ένα
κοινόχρηστο αντικείμενο, το \code{Bid}, άρα περνά από τη συναίνεση
(Ενότητα~\ref{sec:sui-consensus}), της οποίας η καθυστέρηση σε δημόσιο δίκτυο με δεκάδες
γεωγραφικά διάσπαρτους επικυρωτές είναι της τάξης του μισού δευτερολέπτου κατά τη
δημοσιευμένη αξιολόγηση του \en{Mysticeti}~\cite{mysticeti2023}, στην οποία προστίθεται
ο χρόνος του πορτοφολιού και η απόσταση του χρήστη από τον κόμβο. Ο χρήστης της
εφαρμογής θα δει επομένως τα αποτελέσματα των ενεργειών του σε ένα έως δύο δευτερόλεπτα
αντί για το μισό του τοπικού δικτύου, και η συμπεριφορά υπό ανταγωνισμό θα είναι η ίδια
με την Ενότητα~\ref{sec:stress-contention}, αφού την επιβάλλει το ίδιο μοντέλο εκτέλεσης.
Η προβολή αυτή είναι εκτίμηση και όχι μέτρηση· η επαλήθευσή της στο \en{testnet}, που
τρέχει το ίδιο πρωτόκολλο με το \en{mainnet}, είναι το φυσικό επόμενο βήμα.

\section{Αποτελέσματα των σεναρίων}
\label{sec:scenario-results}

Τα δύο εξωτερικά επίπεδα επαλήθευσης της Ενότητας~\ref{sec:verification} εκτελέσθηκαν
στο τοπικό δίκτυο μετά την τελική έκδοση του συμβολαίου και της εφαρμογής. Και τα 20
σενάρια συμβολαίου και οι 17 ροές διεπαφής ολοκληρώθηκαν επιτυχώς, χωρίς σφάλμα
εκτέλεσης της εφαρμογής στην κονσόλα του φυλλομετρητή. Ο
Πίνακας~\ref{tab:scenarios} συνοψίζει τα σενάρια που προσθέτουν γνώση πέρα από τους
ελέγχους μονάδας· τα υπόλοιπα επαναλαμβάνουν στο πραγματικό δίκτυο τις κανονικές ροές
και τους επτά κωδικούς σφαλμάτων.

\begin{table}[!tbp]
\centering
\caption{Σενάρια συμβολαίου στο τοπικό δίκτυο που υπερβαίνουν τους ελέγχους μονάδας}
\label{tab:scenarios}
\small
\begin{tabular}{p{5.6cm}p{8.4cm}}
\toprule
Σενάριο & Παρατηρούμενη συμπεριφορά \\
\midrule
Δύο προσφορές για το ίδιο \en{NFT}· αποδοχή της πρώτης & Η δεύτερη παραμένει στην αλυσίδα. Ο παλιός κάτοχος δεν μπορεί ούτε να την αποδεχθεί ούτε να την απορρίψει, διότι δεν κατέχει πλέον τον στόχο· ο νέος κάτοχος την αποδέχθηκε. \\
Δύο ταυτόχρονες αποδοχές της ίδιας ανοικτής προσφοράς από διαφορετικούς λογαριασμούς & Ακριβώς μία επιτυγχάνει. Η άλλη απορρίπτεται με \code{InputObjectDeleted}, χωρίς καμία μεταβολή κατάστασης πέραν του \en{gas}. \\
Ο στόχος είναι δεσμευμένος σε άλλη προσφορά & Η προσφορά δεν μπορεί να διεκπεραιωθεί από κανέναν, αφού κανείς δεν μπορεί να δώσει τον στόχο ως είσοδο· η ακύρωση από τον προσφέροντα λειτουργεί, και μετά την απελευθέρωση του στόχου η αποδοχή επίσης. \\
Αποδοχή από τον ίδιο τον προσφέροντα & Επιτρέπεται και είναι ακίνδυνη: όλα τα στοιχεία, μαζί με το \en{SUI}, καταλήγουν ξανά σε αυτόν. \\
Προσφορά \code{Bid<MarketplaceNFT, Coin<SUI>>} & Έγινε αποδεκτή με ένα αντικείμενο νομίσματος ως αντάλλαγμα, χωρίς καμία αλλαγή στο συμβόλαιο: οι γενικές παράμετροι δέχονται οποιονδήποτε τύπο \code{key + store}. \\
Λανθασμένα ορίσματα τύπου στην αποδοχή & Απορρίπτονται από την εικονική μηχανή με \code{TypeMismatch} πριν την εκτέλεση. \\
Δεύτερη αποδοχή ήδη καταναλωμένης προσφοράς & Απορρίπτεται κατά τον έλεγχο εισόδων, πριν καν ζητηθεί υπογραφή, επειδή το αντικείμενο \code{Bid} έχει διαγραφεί. \\
\bottomrule
\end{tabular}
\end{table}

Οι ροές διεπαφής ανέδειξαν δύο σημεία που δεν είναι σφάλματα αλλά όρια της τρέχουσας
υλοποίησης. Οι λίστες δεν ανανεώνονται περιοδικά: μια προσφορά που φθάνει ενώ ο χρήστης
βρίσκεται ήδη στη σελίδα «\en{My bids}» εμφανίζεται μόλις ο χρήστης αλλάξει σελίδα ή
επιστρέψει στο παράθυρο. Και όταν ο χρήστης επιχειρεί να αποδεχθεί προσφορά που ο
προσφέρων ακύρωσε στο μεταξύ, η εφαρμογή εμφανίζει το ωμό μήνυμα του κόμβου («\en{Object
\ldots{} has been deleted}») αντί για φιλική διατύπωση, και η παρωχημένη γραμμή παραμένει
μέχρι την επόμενη ανάγνωση. Και τα δύο αντιμετωπίζονται με περιοδική ανανέωση των
ερωτημάτων και καταγράφονται στις μελλοντικές επεκτάσεις.

\section{Κάλυψη απαιτήσεων}
\label{sec:coverage}

Ο Πίνακας~\ref{tab:traceability} αντιστοιχίζει τις απαιτήσεις του
Κεφαλαίου~\ref{ch:apaitiseis} στους ελέγχους που τις επαληθεύουν. Οι αριθμοί στη στήλη
«Έλεγχοι \en{Move}» παραπέμπουν στον Πίνακα~\ref{tab:tests}· οι αριθμοί με πρόθεμα Ν
στον Πίνακα~\ref{tab:tests-nft}. Οι στήλες «\en{SDK}» και
«Διεπαφή» δείχνουν αν η απαίτηση ασκείται από τα σενάρια συμβολαίου μέσω \en{SDK} και από
την αυτοματοποιημένη πλοήγηση, αντίστοιχα (Ενότητα~\ref{sec:verification}).

\begin{table}[H]
\centering
\caption{Αντιστοίχιση απαιτήσεων σε ελέγχους}
\label{tab:traceability}
\small
\begin{tabular}{lp{5.4cm}lcc}
\toprule
Απαίτηση & Περιγραφή & Έλεγχοι \en{Move} & \en{SDK} & Διεπαφή \\
\midrule
ΛΑ1 & Δημιουργία \en{NFT} & Ν1, Ν3 & \yes & \yes \\
ΛΑ2 & Προσφορά με \en{NFTs} ή/και \en{SUI} & 1, 2, 3, 7, 8 & \yes & \yes \\
ΛΑ3 & Στοχευμένη προσφορά & 8, 9, 11, 18 & \yes & \yes \\
ΛΑ4 & Αποδοχή & 6 έως 9, Ν3 & \yes & \yes \\
ΛΑ5 & Απόρριψη & 15, 16, 17 & \yes & \yes \\
ΛΑ6 & Ακύρωση & 13, 14 & \yes & \yes \\
ΛΑ7 έως ΛΑ10 & Προβολές της εφαρμογής & --- & \yes & \yes \\
\midrule
ΜΛΑ1 & Ατομικότητα & 10, 11, 12 & \yes & --- \\
ΜΛΑ2 & Διατήρηση στοιχείων & 1, 7, 13, 15 & \yes & \yes \\
ΜΛΑ3 & Εξουσιοδότηση & 14, 16, 17, 18 & \yes & --- \\
ΜΛΑ5 & Γενικότητα τύπων & 1 έως 18 & \yes & --- \\
ΜΛΑ6 & Ταυτόχρονη πρόσβαση & --- & \yes & --- \\
ΜΛΑ8 & Μία υπογραφή ανά ενέργεια & --- & --- & \yes \\
\bottomrule
\end{tabular}
\end{table}

Τρεις παρατηρήσεις. Η γενικότητα ως προς τους τύπους (ΜΛΑ5) επαληθεύεται από τους
ελέγχους \en{Move}, που χρησιμοποιούν δύο διαφορετικούς τύπους, και από το σενάριο
ανταλλαγής \en{NFT} με νόμισμα· η εφαρμογή επίδειξης διαθέτει έναν μόνο τύπο \en{NFT}, και
οι συλλογές της διαφέρουν ως προς ένα πεδίο. Η ορθότητα υπό ταυτόχρονη πρόσβαση (ΜΛΑ6)
ελέγχεται πειραματικά μόνο στο τοπικό δίκτυο, με δύο λογαριασμούς· κατά τα λοιπά στηρίζεται
στο μοντέλο εκτέλεσης του \en{Sui}, όπως αναλύθηκε στην Ενότητα~\ref{sec:security}. Η
απουσία έμπιστου τρίτου (ΜΛΑ4) και η διαφάνεια (ΜΛΑ7) είναι ιδιότητες του σχεδιασμού που
επαληθεύονται με επιθεώρηση του κώδικα: κανένα από τα δύο \en{modules} δεν περιέχει
συνάρτηση με προνομιακό καλούντα.

\section{Περιορισμοί}
\label{sec:limitations}

\paragraph{Ένας τύπος προσφερομένων ανά προσφορά.} Τα \en{NFTs} που προσφέρονται σε μία
προσφορά πρέπει να είναι του ίδιου τύπου \en{Move}. Προσφορά που συνδυάζει αντικείμενα δύο
διαφορετικών τύπων δεν εκφράζεται με την τρέχουσα υπογραφή.

\paragraph{Η εφαρμογή επίδειξης χρησιμοποιεί έναν τύπο.} Οι τρεις συλλογές της εφαρμογής
είναι τιμές ενός πεδίου του ίδιου τύπου, και η διεπαφή περνά τον \code{MarketplaceNFT} και
στις δύο παραμέτρους. Η ανταλλαγή μεταξύ διαφορετικών τύπων επιδεικνύεται από τους ελέγχους
του συμβολαίου, όχι από τη διεπαφή. Η υποστήριξη αυθαίρετων τύπων στη διεπαφή απαιτεί
ανακάλυψη των \en{NFTs} του χρήστη ανεξαρτήτως τύπου και ανάγνωση του \en{Display} τους.

\paragraph{Ανακάλυψη μέσω συμβάντων.} Κάθε ερώτημα της εφαρμογής διαβάζει όλα τα συμβάντα
ενός τύπου από την αρχή της αλυσίδας και ανακτά όλα τα αντίστοιχα αντικείμενα. Το κόστος
αυξάνει γραμμικά με το συνολικό ιστορικό, όχι με το πλήθος των ενεργών προσφορών. Για μια
επίδειξη είναι επαρκές· για πραγματική χρήση απαιτείται \en{indexer}.

\paragraph{Απουσία λήξης και αντιπρότασης.} Οι προσφορές δεν λήγουν, και ο κάτοχος δεν
μπορεί να απαντήσει με διαφορετικούς όρους· μπορεί μόνο να δεχθεί ή να απορρίψει.

\paragraph{Επιστροφή της εγγύησης αποθήκευσης.} Όπως έδειξε η
Ενότητα~\ref{sec:gas-eval}, η εγγύηση που καταβάλλει ο προσφέρων επιστρέφεται στον
κάτοχο όταν αυτός κλείνει την προσφορά.

\paragraph{Μη επιβολή πολιτικών δημιουργών.} Το πρωτόκολλο μεταβιβάζει αντικείμενα με
\code{public\_transfer} και δεν γνωρίζει για \en{royalties}.

\paragraph{Ανανέωση δεδομένων.} Οι λίστες της εφαρμογής ανανεώνονται μετά από ενέργεια
του ίδιου του χρήστη ή αλλαγή σελίδας, όχι περιοδικά· ενέργειες τρίτων γίνονται ορατές με
καθυστέρηση (Ενότητα~\ref{sec:scenario-results}).

\paragraph{Κάλυψη ελέγχων.} Η ταυτόχρονη πρόσβαση ελέγχεται μόνο σε τοπικό δίκτυο, με
δύο λογαριασμούς· η συμπεριφορά υπό πραγματικό φόρτο δικτύου δεν έχει μετρηθεί. Δεν έχει
γίνει ανεξάρτητος έλεγχος ασφάλειας.

\paragraph{Περιβάλλον αξιολόγησης.} Οι μετρήσεις έγιναν σε τοπικό δίκτυο ενός επικυρωτή,
με παλαιότερη έκδοση πρωτοκόλλου από το \en{mainnet}. Τα ποσά \en{gas} μεταφέρονται με
ακρίβεια (Ενότητα~\ref{sec:mainnet-projection}), οι χρόνοι και ο ρυθμός διεκπεραίωσης
όχι· η εκτίμησή τους για δημόσιο δίκτυο είναι ποιοτική έως ότου επαναληφθούν στο
\en{testnet}.
